\section{Velocity}

\subsection{Why position is not enough}

Knowing the position of a moving point at every time tells us where the point is. But we often want a more local description: how the position is changing at a given moment.

If the position changes a lot during a short time interval, we say that the motion is fast. If the position changes only a little during the same time interval, we say that the motion is slow. If the position changes in a different direction, the motion has a different direction.

Velocity is the mathematical object that describes this change of position.

\subsection{Average velocity}

Let the position of a moving point be \(\mathbf{r}(t)\). Suppose we compare two moments of time, \(t_1\) and \(t_2\), with \(t_1\neq t_2\). The displacement over this interval is
\[
\Delta \mathbf{r}=\mathbf{r}(t_2)-\mathbf{r}(t_1),
\]
and the elapsed time is
\[
\Delta t=t_2-t_1.
\]

\begin{definitionbox}
The average velocity of a moving point on the interval \([t_1,t_2]\), with \(t_1\neq t_2\), is
\[
\mathbf{v}_{\mathrm{avg}}
=
\frac{\Delta \mathbf{r}}{\Delta t}
=
\frac{\mathbf{r}(t_2)-\mathbf{r}(t_1)}{t_2-t_1}.
\]
\end{definitionbox}

\begin{figure}[htbp]
\centering
\begin{tikzpicture}[scale=1.19, transform shape, >={Stealth[length=4.5mm,width=3mm]}]
    \draw[step=1, gray!25, very thin] (-3.6,-0.5) grid (3.6,3.6);
    \draw[->] (-3.7,0) -- (3.9,0) node[right] {$x$};
    \draw[->] (0,-0.5) -- (0,3.9) node[above] {$y$};
    \draw[domain=0:180, smooth, samples=100] plot ({3*cos(\x)},{3*sin(\x)});
    \coordinate (O) at (0,0);
    \coordinate (Pone) at ({3*cos(45)},{3*sin(45)});
    \coordinate (Ptwo) at ({3*cos(120)},{3*sin(120)});
    \draw[-{Stealth[length=4.5mm,width=3mm]}, line width=0.9pt] (O) -- (Pone);
    \draw[-{Stealth[length=4.5mm,width=3mm]}, line width=0.9pt] (O) -- (Ptwo);
    \draw[-{Stealth[length=5.2mm,width=3.5mm]}, line width=1.2pt] (Pone) -- (Ptwo);
    \fill (Pone) circle (2pt);
    \fill (Ptwo) circle (2pt);
    \node[below right] at (1.05,1.05) {$\mathbf{r}(t_1)$};
    \node[left] at (-0.95,1.55) {$\mathbf{r}(t_2)$};
    \node[above] at (0.35,2.65) {$\Delta\mathbf{r}$};
    \node[below left] at (0,0) {$O$};
\end{tikzpicture}

\caption{Displacement between two positions of a moving point.}
\label{fig:kinematics-displacement-between-positions}
\end{figure}

The average velocity is a vector. Its direction is the direction of the displacement \(\Delta\mathbf{r}\), and its magnitude tells us how much displacement occurs per unit time during the interval. In one dimension,
\[
v_{\mathrm{avg}}=
\frac{x(t_2)-x(t_1)}{t_2-t_1}.
\]

\begin{remarkbox}
Average velocity depends on the interval. A large interval describes overall change; a smaller interval gives a more local description.
\end{remarkbox}

\subsection{Instantaneous velocity}

Average velocity describes motion over a time interval. To describe motion at one instant, we shrink the interval.

\begin{definitionbox}
Let \(\mathbf{r}\) be differentiable at time \(t\). The instantaneous velocity at time \(t\) is
\[
\mathbf{v}(t)
=
\lim_{\Delta t\to 0}
\frac{\mathbf{r}(t+\Delta t)-\mathbf{r}(t)}{\Delta t}
=
\frac{d\mathbf{r}}{dt}(t).
\]
\end{definitionbox}

The limiting process is essential. We are not dividing by zero. We first compute the average velocity over a nonzero time interval \(\Delta t\), and only then study what this expression approaches as \(\Delta t\) becomes arbitrarily small.

If
\[
\mathbf{r}(t)=(x(t),y(t),z(t)),
\]
then
\[
\mathbf{v}(t)=\left(\frac{dx}{dt},\frac{dy}{dt},\frac{dz}{dt}\right).
\]
In two dimensions,
\[
\mathbf{v}(t)=\left(\frac{dx}{dt},\frac{dy}{dt}\right),
\]
and in one dimension,
\[
v(t)=\frac{dx}{dt}.
\]

Velocity is therefore the derivative of position with respect to time.

\subsection{Velocity as a vector attached to the moving point}

Velocity is a vector quantity. It has magnitude and direction. At each time \(t\), the point has a position \(\mathbf{r}(t)\) and a velocity \(\mathbf{v}(t)\). It is useful to imagine the velocity vector as attached to the point at its current position.

The velocity vector is tangent to the trajectory. The trajectory tells us the path. The velocity tells us the local direction in which the point is moving along that path.

\begin{remarkbox}
For motion along a circle, velocity is tangent to the circle, not directed toward the center. The direction of velocity changes from point to point, even if the speed remains constant.
\end{remarkbox}

\subsection{Speed}

\begin{definitionbox}
The speed of a moving point at time \(t\) is the magnitude of its velocity vector:
\[
v(t)=|\mathbf{v}(t)|.
\]
\end{definitionbox}

If
\[
\mathbf{v}(t)=(v_x(t),v_y(t),v_z(t)),
\]
then
\[
|\mathbf{v}(t)|=\sqrt{v_x^2(t)+v_y^2(t)+v_z^2(t)}.
\]
Speed is not a vector. It is a non-negative number. It tells us how fast the point is moving, but not in which direction. Velocity tells us both: how fast and in which direction.

\subsection{Examples}

\begin{examplebox}
For one-dimensional motion
\[
x(t)=x_0+at,
\]
where \(x_0\) and \(a\) are constants, the velocity is
\[
v(t)=\frac{dx}{dt}=a.
\]
If \(x(t)=1+2t\), then \(v(t)=2\). The number \(2\) is the rate at which position changes with respect to time.
\end{examplebox}

\begin{examplebox}
For
\[
\mathbf{r}(t)=(t,t^2),
\]
we have
\[
\mathbf{v}(t)=\left(\frac{dx}{dt},\frac{dy}{dt}\right)=(1,2t).
\]
Thus
\[
\mathbf{v}(0)=(1,0),\qquad \mathbf{v}(1)=(1,2),\qquad \mathbf{v}(2)=(1,4).
\]
The horizontal component is constant, while the vertical component changes.
\end{examplebox}

\begin{examplebox}
For circular motion
\[
\mathbf{r}(t)=(R\cos t,R\sin t),
\]
the velocity is
\[
\mathbf{v}(t)=\frac{d\mathbf{r}}{dt}=(-R\sin t,R\cos t).
\]
At \(R=2\) and \(t=\pi/3\),
\[
\mathbf{r}\left(\frac{\pi}{3}\right)=(1,\sqrt{3}),
\qquad
\mathbf{v}\left(\frac{\pi}{3}\right)=(-\sqrt{3},1).
\]
Their dot product is
\[
(1,\sqrt{3})\cdot(-\sqrt{3},1)=0,
\]
so at this point the velocity vector is perpendicular to the position vector.
\end{examplebox}

\begin{figure}[htbp]
\centering
\begin{tikzpicture}[scale=1.61, transform shape, >={Stealth[length=4.5mm,width=3mm]}]
    \draw[step=1, gray!25, very thin] (-2.6,-0.5) grid (2.6,3.2);
    \draw[->] (-2.7,0) -- (2.9,0) node[right] {$x$};
    \draw[->] (0,-0.5) -- (0,3.4) node[above] {$y$};
    \draw (0,0) circle (2);
    \coordinate (O) at (0,0);
    \coordinate (P) at (1,{sqrt(3)});
    \coordinate (Vend) at ({1-sqrt(3)},{sqrt(3)+1});
    \draw[-{Stealth[length=4.5mm,width=3mm]}, line width=0.9pt] (O) -- (P);
    \draw[-{Stealth[length=5.2mm,width=3.5mm]}, line width=1.2pt] (P) -- (Vend);
    \fill (P) circle (2pt);
    \node[above right] at (P) {$P=(1,\sqrt{3})$};
    \node[below right] at (0.55,0.75) {$\mathbf{r}$};
    \node[above left] at (0.25,2.35) {$\mathbf{v}$};
    \node[below left] at (O) {$O$};
\end{tikzpicture}

\caption{Velocity tangent to a circular trajectory.}
\label{fig:kinematics-circular-velocity-tangent}
\end{figure}

The speed in this circular motion is
\[
|\mathbf{v}(t)|=\sqrt{(-R\sin t)^2+(R\cos t)^2}=R.
\]
The speed is constant, although the velocity vector changes direction.
